\label{sec:tokens}
This section develops the simplest framing of the problem: every input slot is a vertex, the weights are weighted
edges, a layer is a quadratic form on the vertex graph, the next layer is read from the two-token (Kikuchi) lift of
that graph, and a second tier---the input law---sets the environment in which the lower tier's walks take place. The
framing is exact: it is the Wiener chaos (bosonic Fock) structure of Gaussian space, with tokens as Hermite quanta.
It reorganises everything above in a few lines, and it locates the obstruction precisely: it is a contact of tokens at
a hidden neuron.

\subsection{Dictionary}
Let $V=[n]$ be the input slots. By the Wiener--It\^o isomorphism
\[
L^2(\R^n,\gamma)\;\cong\;\bigoplus_{k\ge0}\Sym^k\big(\ell^2(V)\big),\qquad
F=\sum_k\big\langle f_k,\,{:}x^{\otimes k}{:}\big\rangle,\qquad f_k=\tfrac1{k!}\,\E\big[\nabla^kF\big]\ \ \text{(Stroock)},
\]
the $k$-th chaos is spanned by the multivariate Hermite polynomials $\He_\alpha$, $|\alpha|=k$: configurations of $k$
indistinguishable tokens on the vertices, repetitions allowed (the bosonic $k$-token Kikuchi level). The Gaussian input
is the vacuum $\Omega=1$; $a_i=\partial_i$ removes a token at $i$, $a_i^\dagger=x_i-\partial_i$ adds one, $N=\sum_ia_i^\dagger a_i$ is the
Ornstein--Uhlenbeck number operator. A weighted graph on $V$ is a function on the pair groupoid $V\times V$, i.e.\ an element
of $M_n=\mathbb C[V\times V]$; it moves one token along an edge as $d\Gamma(A)$ and all tokens as $\Gamma(A)$. Two-token kernels
are quadratic forms, i.e.\ elements of the groupoid algebra, and contracting two kernels in $r$ shared token slots,
$f\otimes_rg$, is groupoid convolution ($Q\otimes_1Q'=QQ'$, $Q\otimes_2Q'=\Tr QQ'$). Positive definite functions on the pair
groupoid act by Schur multiplication as completely positive maps; the unit space (the diagonal) is where two tokens sit
on one vertex.

\subsection{Layer one: each neuron is a one-mode token tower, and its mean is a trace}
\begin{proposition}[Ridge towers]\label{prop:ridge}
With $c_k=\E[\relu(Z)\He_k(Z)]$, i.e.\ $(c_0,c_1,c_2,c_3,c_4,c_5,c_6,\dots)=(\varphi_0,\tfrac12,\varphi_0,0,-\varphi_0,0,3\varphi_0,\dots)$,
$\varphi_0=1/\sqrt{2\pi}$, $c_k=\He_{k-2}(0)\varphi_0$ for $k\ge2$,
\[
h_{1,a}=|w_a|\sum_k\frac{c_k}{k!}\big(a^\dagger(\hat w_a)\big)^k\Omega ,\qquad
f_2(h_{1,a})=Q_a=\frac{w_aw_a^{\top}}{2\sqrt{2\pi}\,|w_a|}.
\]
Each layer-one neuron lives in the one-mode Fock space of its own edge vector, and its quadratic form is the rank-one
graph of that edge. Its degree-$k$ coefficient is $\E[\relu^{(k)}]$: the gate for $k=1$, the kink density for $k=2$, and
derivatives of the kink for $k\ge3$.
\end{proposition}
\begin{proof}
$\relu(|w|s)=|w|\relu(s)$ with $s=\hat w\cdot x\sim N(0,1)$, the one-dimensional Hermite expansion, and $\He_k(\hat w\cdot x)=
(a^\dagger(\hat w))^k\Omega$. $c_k=\E[\relu^{(k)}]=\E[\delta^{(k-2)}]=\He_{k-2}(0)\varphi_0$ by Gaussian integration by parts.
\end{proof}

\begin{theorem}[Euler--Fock ladder]\label{thm:ladder}
If $F$ is positively homogeneous of degree one then $(N+\sum_ia_ia_i)F=F$, i.e.
\[
f_0=2\tr f_2,\qquad \tr_{23}f_3=0,\qquad f_k=-\frac{(k+1)(k+2)}{k-1}\,\tr_{k+1,k+2}f_{k+2}\quad(k\ge2).
\]
In particular the mean of every neuron of every layer is twice the trace of its quadratic form: the mean is a closed
two-token loop, the two tokens of the quadratic form placed on the same input vertex and summed.
\end{theorem}
\begin{proof}
$x\cdot\nabla=\sum_i(a_i^\dagger+a_i)a_i=N+\sum_ia_ia_i$, and Euler's relation $x\cdot\nabla F=F$ holds for one-homogeneous $F$.
$\sum_ia_ia_i$ maps the chaos-$(k+2)$ component with kernel $f_{k+2}$ to chaos $k$ with kernel $(k+2)(k+1)\tr f_{k+2}$, the trace over
two token slots. Compare the chaos-$k$ components. For $h_{1,a}$: $2\tr Q_a=|w_a|/\sqrt{2\pi}$, and $-12\tr f_4=f_2$ with
$f_4=-\tfrac{\varphi_0}{24}|w_a|\hat w_a^{\otimes4}$.
\end{proof}

\subsection{Layer two: the quadratic form is the graph of two-step walks}
\begin{theorem}[Layer-two tokens]\label{thm:layer2}
For $z_{2,c}=\sum_aW_2(c,a)h_{1,a}$, with $\omega=W_2(c,\cdot)^{\top}$,
\begin{gather*}
z_{2,c}=m_c+\langle f_c,x\rangle+{:}x^{\top}Q_cx{:}+(\text{chaos}\ge4),\\
f_c=\tfrac12W_1^{\top}\omega,\qquad
Q_c=W_1^{\top}D_cW_1,\qquad D_c=\diag\Big(\frac{W_2(c,a)}{2\sqrt{2\pi}|w_a|}\Big).
\end{gather*}
$Q_c(i,j)=\sum_aW_1(a,i)\,D_c(a)\,W_1(a,j)$ is the weighted graph of two-step walks $i\to a\to j$ through the hidden layer,
turning at $a$ with weight $D_c(a)$. By Sylvester's identity every spectral quantity of $Q_c$ is one of $D_cG$, where
$G=W_1W_1^{\top}$ is the codegree graph of the hidden neurons (overlaps of their input supports; normalised, the cosines
$\rho_{ab}$).
\end{theorem}
\begin{proof}
Linearity of the chaos decomposition and Proposition~\ref{prop:ridge} ($\E g_{1,a}=\tfrac12$).
\end{proof}

\begin{theorem}[Cumulants are walk sums]\label{thm:walks}
(i) For $z=m+\langle f,x\rangle+x^{\top}Qx-\tr Q$,
\[
\log\E e^{itz}=it\,(m-\tr Q)-\tfrac12\log\det(I-2itQ)-\tfrac{t^2}2\,f^{\top}(I-2itQ)^{-1}f ,
\]
\begin{align*}
\kap_r(z)&=2^{r-1}(r-1)!\,\Tr Q^r+r!\,2^{r-3}f^{\top}Q^{r-2}f\\
&=2^{r-1}(r-1)!\,\Tr(D_cG)^r+r!\,2^{r-5}\,\omega^{\top}G(D_cG)^{r-2}G\omega :
\end{align*}
closed walks on the hidden codegree graph, and open walks whose two ends are the linear token.
(ii) Exactly (all chaoses), $\kap_r(z_{2,c})=\sum_{a_1,\dots,a_r}\prod_j\omega_{a_j}|w_{a_j}|\;\kap\big(\relu(s_{a_1}),\dots,\relu(s_{a_r})\big)$, and for
distinct indices the joint cumulant of the ridge functions is the sum over connected multigraphs on $\{a_1,\dots,a_r\}$
with weight $\rho^m/m!$ on an edge of multiplicity $m$ and weight $c_{\deg(v)}$ on a vertex of degree $\deg(v)$ (the diagram
formula). For $r=3$ the leading graphs are the three paths, centred at a hub of degree two (a kink, $c_2=\varphi_0$) with
two gated arms ($c_1=\tfrac12$), and the triangle ($c_2^3$). Repeated indices (two or three tokens on the same hidden
neuron: contacts) are given by the exact single-neuron and pair laws, not by the quadratic truncation.
\end{theorem}
\begin{proof}
(i) Diagonalise $Q$; each mode is a shifted noncentral $\chi^2$, and completing the square gives the displayed
log-characteristic function. Expand $\log\det(I-2itQ)=-\sum_r(2it)^r\Tr Q^r/r$ and $(I-2itQ)^{-1}=\sum_s(2it)^sQ^s$, and read off $r!$
times the coefficient of $(it)^r$. Substitute $Q=W_1^{\top}D_cW_1$ and $f=\tfrac12W_1^{\top}\omega$. (ii) Mehler's formula for
Gaussian vectors with unit variances: $\E\prod_jg_j(s_j)=\sum\prod_{v}\E[g_v^{(\deg v)}]\prod_{e}\rho_e^{m_e}/m_e!$ over multigraphs;
cumulants keep the connected ones.
\end{proof}
\begin{center}\small
\begin{tabular}{lccccc}\toprule
width & $\kap_3(z_{2,c})$ (MC) & self & contacts & distinct neurons & leading diagrams\\\midrule
$24$ & $-0.521$ & $-0.465$ & $-0.115$ & $0.059$ & $0.064$\\
$96$ & $0.210$ & $0.078$ & $0.121$ & $0.0115$ & $0.0133$\\\bottomrule
\end{tabular}
\end{center}
(Self: all three indices on one neuron; contacts: two of three on one neuron, by the exact pair law; distinct: Monte
Carlo minus the two exact parts; leading diagrams: hub paths and triangles.) The distinct-neuron part is the walk sum, to the accuracy of its leading diagrams; at these widths the contacts carry
most of a single pre-activation's third cumulant. The quadratic truncation of (i) applied to the whole of $z_{2,c}$ gives
$0.070$ against the true $-0.521$ ($n=24$): the walk picture is right for exchange between distinct neurons and wrong for
contacts, which must be computed exactly.

\subsection{How many tokens: depth heats the tower}
\begin{proposition}[Mean token number]\label{prop:tokens}
For $F\in L^2(\gamma)$ the mean token number of its centred part is $\bar N(F)=\E|\nabla F|^2/\Var F$ (since
$\E|\nabla F|^2=\sum_kk\,k!\|f_k\|^2$ and $\Var F=\sum_{k\ge1}k!\|f_k\|^2$). In the infinite-width limit the pre-activations of layer $l$
have
\[
\bar N_l=\frac1{1-\cos\theta_{l-1}},\qquad \theta_0=\tfrac\pi2,\quad \cos\theta_{k+1}=J(\theta_k)/\pi ,
\qquad \bar N_l\simeq\frac{2l^2}{9\pi^2}\ (l\to\infty),
\]
the pair temperature of Proposition~\ref{prop:heattrace}.
\end{proposition}
\begin{proof}
For a one-homogeneous random field with kernel $K(x,y)=\tfrac{|x||y|}nk(\hat x\cdot\hat y)$ differentiating at $y=x$, $|x|^2=n$, gives
$\E|\nabla_xz|^2=\sum_i\partial_{x_i}\partial_{y_i}K|_{y=x}=\tfrac1n\big(k(1)+(n-1)k'(1)\big)$ while $\E z^2=k(1)$; the normalised dual kernel of \textsc{ReLU} has derivative $1$ at $t=1$ at every layer (criticality), so
$\E|\nabla z|^2=\E z^2(1+O(1/n))$. The squared mean $m^2$ is the kernel at two independent inputs, whose cosine tends to $0$, so
$m^2/\E z^2\to\cos\theta_{l-1}$ with the angle recursion of Theorem~\ref{thm:folding} started at $\pi/2$. Then
$\bar N=\E z^2/(\E z^2-m^2)$; the asymptotics follow from $\beta_l\simeq9\pi^2/(2l^2)$.
\end{proof}
\begin{center}\small
\begin{tabular}{lccccccc}\toprule
layer & $1$ & $2$ & $3$ & $5$ & $9$ & $13$ & $16$\\\midrule
infinite width & $1.00$ & $1.47$ & $1.98$ & $3.13$ & $6.04$ & $9.72$ & $12.99$\\
width $256$ (Monte Carlo) & $1.00$ & $1.47$ & $2.01$ & $3.12$ & $6.26$ & $9.46$ & $10.56$\\\bottomrule
\end{tabular}
\end{center}
So the token tower does not truncate: the last pre-activation of a sixteen-layer network carries eleven to thirteen
tokens on average, far from the two of a quadratic form (and far below the $2^{l-1}$ of naive degree doubling).
Its marginal law is nevertheless close to Gaussian (skewness about $0.1$ on network $0$). The two facts are reconciled by the
fourth-moment theorem of Nualart and Peccati \cite{NP}: within a chaos, Gaussianity is controlled by the contractions $f\otimes_rf$
($1\le r<k$), the overlaps of a kernel's token configurations with themselves, i.e.\ token codegrees, and not by the token
number. The Gaussian closure works because these contractions are small at He initialisation, and it fails exactly
where they are not: at contacts, and along the collective mode.

\subsection{Two tiers: the input law is the environment}
\begin{theorem}[One-token plus two-token]\label{thm:tiers}
For a one-homogeneous $F$ and every Gaussian input law $N(m,\Sigma)$,
\[
\E_{N(m,\Sigma)}F=\big\langle m,\ \E\nabla F\big\rangle+\big\langle\Sigma,\ \E\nabla^2F\big\rangle ,
\]
with the distributional Hessian (walls). Equivalently, the exact mean $u(m,\Sigma)$ satisfies Euler's relation and the heat
equation $\partial_{\Sigma_{ij}}u=\tfrac12\partial_{m_i}\partial_{m_j}u$; in token language, adding a token pair at $(i,j)$ to the input
state has the same effect as two successive displacements along $i$ and $j$.
\end{theorem}
\begin{proof}
$\E F=\E[x\cdot\nabla F]=\langle m,\E\nabla F\rangle+\E[(x-m)\cdot\nabla F]$ and Stein's identity for $N(m,\Sigma)$,
$\E[(x-m)g]=\Sigma\,\E\nabla g$, with $g=\nabla F$. The heat equation is $\partial_\Sigma\E F(m+\Sigma^{1/2}Z)=\tfrac12\E\nabla^2F$.
\end{proof}

The upper tier is the input law; it sets the environment of the lower tier, namely the kink densities and the
gate-pair kernels of every layer. A moment chain $e(m,\Sigma)$ obeys Euler's relation (scale covariance) but not the heat
equation: its pair response and its double-displacement response disagree, and that disagreement is the heat defect of
stage 9. Along a localization (Eldan's balls of radius $r$ about a moving centre, Theorem~\ref{thm:uniform}) the environment
cools. Gates freeze, the contact weights $\Phi(1-\Phi)$ vanish away from walls, the two-token term transfers into the
one-token term, and at $r=0$ the walk is free and every chain is exact. The Euler--Stein correction integrates the
chain's mismatch over this cooling, and its constant $1/(d+2)$ records the token order of the mismatch at the kink.

\subsection{The two-token walk and its environment}
Pathwise, the gradient Gram $\Gamma_l=J_lJ_l^{\top}$ ($J_l=\partial z_l/\partial x$) obeys $\Gamma_l=W_lD_{l-1}\Gamma_{l-1}D_{l-1}W_l^{\top}$: two tokens
walk up the layered graph, each picking up weights and gates. In a factorised environment its expectation is the
completely positive map $\Gamma\mapsto W(K\circ\Gamma)W^{\top}$, where the gate-pair kernel $K(a,b)=\mathbb P(z_a>0,z_b>0)$ is a positive definite
function on the pair groupoid. The tetrachoric (Mehler) series splits it into three interactions:
\[
K=\underbrace{\Phi\Phi^{\top}}_{\text{independent walkers}}+\underbrace{\diag(\Phi-\Phi^2)}_{\text{contact}}+\underbrace{\sum_{k\ge1}\frac{1}{k!}\big(\rho^{\circ k}\circ\psi_k\psi_k^{\top}\big)_{\rm off}}_{\text{exchange of $k$ quanta}},
\qquad\psi_k=\varphi(\alpha)\He_{k-1}(\alpha).
\]
The first term is a congruence (the two tokens move independently with mean gates). The second acts only when both
tokens sit on one neuron, and it is diagonal, so it is sparse in the two-token graph. The third is dense, with
off-diagonal entries of order $n^{-k/2}$. The mean reads the unit space of the final quadratic form (its trace).

\begin{proposition}[Independent walkers are wrong at order one, by exactly the contact term]\label{prop:contactterm}
For two layers, let $u^{\rm walk}_c=\Phi_cm_c+f_c(0)\,\E|\nabla z_{2,c}|^2$ be the kink current with independent walkers, and
$u^{\rm GC}_c=\Phi_cm_c+f_c(0)\sigma_c^2$ the Gaussian closure (same exact $m_c,\sigma_c$, $\alpha_c=m_c/\sigma_c$). Then
\[
u^{\rm walk}_c-u^{\rm GC}_c=f_c(0)\big(\E|\nabla z_{2,c}|^2-\sigma_c^2\big)
=f_c(0)\sum_aW_2(c,a)\,\mu_a\big(m_c-\E[z_{2,c}\mid z_{1,a}=0]\big),
\]
whose contact part ($b=a$ in $z_{2,c}=\sum_bW_2(c,b)h_{1,b}$) is $f_c(0)\sum_aW_2(c,a)^2\mu_a^2=O(1)$, while the exchange part
($b\neq a$) is second order in $\rho_{ab}$.
\end{proposition}
\begin{proof}
Stein's identity gives $\E z^2=\E|\nabla z|^2+\E[z\Delta z]$ for the one-homogeneous $z=z_{2,c}$, and
$\Delta z=\sum_aW_2(c,a)\delta(z_{1,a})|w_a|^2$ with $f_{1,a}(0)|w_a|^2=\mu_a$; hence $\E|\nabla z|^2-\sigma^2=m^2-\sum_aW_2(c,a)\mu_a\E[z\mid z_{1,a}=0]$, and
$m=\sum_aW_2(c,a)\mu_a$. Conditioning on $z_{1,a}=0$ sets $h_{1,a}=0$ (the contact) and shifts the other neurons through their
correlations with $z_{1,a}$, which at zero mean changes their means at order $\rho^2$.
\end{proof}
Numerically (width $256$, $2^{20}$ inputs): the Gaussian closure has rms error $1.7\cdot10^{-3}$, independent walkers
$1.9\cdot10^{-1}$; $\E|\nabla z|^2-\sigma^2$ and the contact sum agree in mean ($0.639$ and $0.639$, correlation $0.999$); removing the contact term
brings the walkers to within $7.8\cdot10^{-4}$ of the Gaussian closure. So the Gaussian closure is the two-token walk with
exact contacts. Re-projecting each layer onto a fresh one-token (Gaussian) vector with the same covariance is precisely
what handles self-pinning, and the third-cumulant chain adds first-order exchange.

\subsection{What sparsifies, and what does not}
\begin{enumerate}[leftmargin=*]
\item \textbf{Two-token objects close.} Every sum
$\sum_{l'}T_{l\leftarrow l'}X_{l'}T_{l\leftarrow l'}^{\top}$ accumulates into one $n\times n$ recursion ($n^3$ per layer): the covariance, the gradient Gram, the Gaussian closure, and any hub sum whose hub coefficient is known
when the hub is born.
\item \textbf{Layer two is cheap.} The open walk of length two (the star of Theorem~\ref{thm:walks}) is
\[\textstyle\sum_aD_c(a)\,(W_2G)(c,a)^2 ;\] its hub coefficient is $W_2$ itself, known at birth, so all $n$ of them cost one product
$W_2G$ and a Hadamard square.
\item \textbf{From layer three on the hub coefficient is a walk.} A hub born at layer $l'$ and read at layer $l$ needs
three tokens that meet at the hub, each transported from $l'$ to $l$ (two arms and the reader). Nothing accumulates, and
this is the third-cumulant memory. The trace over input slots (the operation the Gaussian input performs) contracts
input vertices, never hidden hubs, so no operation on the input removes it.
\end{enumerate}
\begin{center}\small
\begin{tabular}{llc}\toprule
token-walk level & estimator & network $0$, MSE\\\midrule
independent walkers & factorised kink current & $O(1)$ (two layers)\\
$+$ exact contacts & Gaussian closure & $4.1\cdot10^{-6}$\\
$+$ first-order exchange & third-cumulant chain (radial $\kap_4$ off/on) & $5.6\cdot10^{-7}$ / $3.5\cdot10^{-8}$\\
$+$ cooling, $d=2$ & Euler--Stein merge, $a\approx\tfrac14$ & $6.6\cdot10^{-9}$\\\bottomrule
\end{tabular}
\end{center}
